\section*{Hoofdstuk \ref{ch:b3:membrane-traffic} --- Membraanverkeer en het sorteren van eiwitten}

\begin{solution}{exo:b3:membrane-traffic:1}
Hydrofoob stuk aan de aminoterminus: \omterm{def:b3:membrane-traffic:signals}{signaalpeptide}, het endoplasmatisch
reticulum in (en standaard verder naar het oppervlak). PKKKRKV:
\omterm{def:b3:membrane-traffic:signals}{kernlokalisatiesignaal}, de kern in door de porie. Amfipathische
aminoterminale helix: mitochondriale \omterm{def:b3:membrane-traffic:signals}{presequentie}, de mitochondriale
matrix in. KDEL: terughaling uit het golgiapparaat naar het reticulum.
Mannose-6-fosfaat: van het trans-golginetwerk naar het \omterm{def:b3:membrane-traffic:endocytosis}{endosoom} en het
\omterm{def:b3:membrane-traffic:endocytosis}{lysosoom}.
\end{solution}

\begin{solution}{exo:b3:membrane-traffic:2}
(1) Zonder membranen was de vertaalde keten ongeveer twintig residuen
langer dan het uitgescheiden eiwit: er wordt tijdens de secretie een
segment verwijderd. (2) Met membranen aanwezig tijdens de translatie had
het product de rijpe grootte en was het beschermd tegen protease: het
was de blaasjes binnengegaan en had het segment daarbinnen verloren.
(3) Membranen die na de translatie werden toegevoegd deden geen van
beide: de intrede is aan de synthese gekoppeld --- cotranslationeel ---
en het segment, het \omterm{def:b3:membrane-traffic:signals}{signaalpeptide}, werkt alleen op een ontluikende
keten.
\end{solution}

\begin{solution}{exo:b3:membrane-traffic:3}
Reticulum naar golgiapparaat: \omterm{def:b3:membrane-traffic:coats}{COPII}, voorwaarts. Golgiapparaat naar
reticulum: \omterm{def:b3:membrane-traffic:coats}{COPI}, retrograad. Trans-golginetwerk naar \omterm{def:b3:membrane-traffic:endocytosis}{endosoom}:
\omterm{def:b3:membrane-traffic:coats}{clathrine} met zijn adaptoren, voorwaarts naar het \omterm{def:b3:membrane-traffic:endocytosis}{lysosoom}.
Plasmamembraan naar \omterm{def:b3:membrane-traffic:endocytosis}{endosoom}: \omterm{def:b3:membrane-traffic:coats}{clathrine}, naar binnen (endocytose).
\end{solution}

\begin{solution}{exo:b3:membrane-traffic:4}
Suikers worden alleen binnen het lumen van het reticulum en het
golgiapparaat toegevoegd. De lumenzijde van elk blaasje blijft door het
knoppen en versmelten heen de lumenzijde, en wanneer een blaasje met het
plasmamembraan versmelt wordt zijn lumenzijde de buitenkant van de cel;
de richting wordt nooit omgekeerd, zodat wat binnen is geglycosyleerd
buiten belandt.
\end{solution}

\begin{solution}{exo:b3:membrane-traffic:5}
$t^{*} = \ln(0.05/0.2)/(0.05 - 0.2) = \ln 4/0.15 = \qty{9.2}{min}$;
$B(t^{*}) = (0.2/(-0.15))(e^{-1.85} - e^{-0.46}) = 0.63$. Na \qty{60}{min}
is $A \approx 0$, $B = 1.33\,(e^{-3} - e^{-12}) = 0.066$, dus $C = 0.93$.
\end{solution}

\begin{solution}{exo:b3:membrane-traffic:6}
Het \omterm{def:b3:membrane-traffic:signals}{signaalpeptide} zet de aminoterminus in het lumen; elk hydrofoob
stuk is afwisselend een stoptransfer of een starttransfer, zodat er
vier transmembraanhelices zijn (ongeveer de residuen 30--52,
90--112, 150--172, 210--232). Segmenten: aminoterminus lumenaal;
52--90 cytosolisch; 112--150 lumenaal; 172--210 cytosolisch;
carboxyterminus lumenaal. Glycosylering alleen op de lumenale segmenten
--- de aminoterminus, 112--150 en de carboxyterminus --- die aan het
oppervlak extracellulair worden.
\end{solution}

\begin{solution}{exo:b3:membrane-traffic:7}
$J_{\max} = \num{30000}\times 0.25\times 0.05/0.30 = 1250$ per minuut;
deel aan het oppervlak $k_{r}/(k_{i} + k_{r}) = 0.05/0.30 = 0.17$. $R$
verdubbelen verdubbelt $J$ precies; de terugkeersnelheid is de
flessenhals (de meeste receptoren zitten binnen), en $k_{r}$ verdubbelen
zou $\num{30000}\times 0.25
\times 0.1/0.35 \approx 2140$ geven, een factor $1.7$.
\end{solution}

\begin{solution}{exo:b3:membrane-traffic:8}
(a) Geen receptor: de gemerkte hydrolasen volgen de standaardroute en
worden uitgescheiden; de lysosomen missen enzymen en lopen vol
onverteerd materiaal --- een stapelingsziekte. (b) Geen
fosfotransferase: hetzelfde fenotype door een ander gebrek (de
I-celziekte), met de enzymen ongemerkt en uitgescheiden. (c)
Geneutraliseerde endosomen: de receptor kan zijn vracht niet loslaten,
wordt dus niet gerecycled en de hydrolasen worden verkeerd afgeleverd
of uitgescheiden; ook LDL- en andere receptoren houden op met rondgaan.
\end{solution}

\begin{solution}{exo:b3:membrane-traffic:9}
Het tyrosine hoort bij het staartmotief dat de clathrine-adaptor
herkent; zonder dat \omterm{def:b3:bioinformatics:motif}{motief} wordt de receptor niet in beklede putjes
verzameld. De receptoren binden LDL normaal maar liggen gelijkmatig
over het oppervlak verspreid, buiten de putjes gehouden, en het ligand
blijft buiten --- een ziekte van het sorteren, niet van het binden.
\end{solution}

\begin{solution}{exo:b3:membrane-traffic:10}
Met $k_{1} = k_{2} = k$ geeft de poging $B = kt\,e^{-kt}$:
$\mathrm{d}B/\mathrm{d}t
= k e^{-kt} - k^{2}t e^{-kt} = kA - kB$, zoals vereist, met $B(0) = 0$.
Maximum waar $1 - kt = 0$: $t^{*} = 1/k$, $B = e^{-1} = 0.37$. In het
algemeen is
$C(t) = 1 - (k_{2}e^{-k_{1}t} - k_{1}e^{-k_{2}t})/(k_{2} - k_{1})$,
wat symmetrisch is onder verwisseling van $k_{1}$ en $k_{2}$: de
korrelcurve alleen kan niet zeggen welk van de twee compartimenten het
snelle is.
\end{solution}

\begin{solution}{exo:b3:membrane-traffic:11}
Elk blaasje heeft oppervlak $4\pi(50\,\text{nm})^{2} = \qty{0.031}{\micro
m^2}$; duizend per minuut is \qty{31}{\micro m^2/min}, dus het hele
oppervlak in $3000/31 \approx \qty{95}{min}$. De cel krimpt niet omdat
er door exocytose (recyclingblaasjes) even snel membraan terugkomt ---
een evenwicht van stromen. Elk blaasje bevat $5\times 10^{-19}$ L;
duizend per minuut gedurende een uur is $3\times 10^{-14}$ L, ongeveer
\qty{1.5}{\%} van het volume van een cel van \qty{2}{pL} per uur.
\end{solution}

\begin{solution}{exo:b3:membrane-traffic:12}
Botulinetoxine werkt in het motorische uiteinde: er komt geen
acetylcholine vrij, de spier krijgt geen opdracht en blijft slap ---
een slappe verlamming. Tetanustoxine wordt langs het motorische axon
omhoog gedragen en over naar de remmende schakelneuronen die de
motorische neuronen normaal in toom houden; met hun \omterm{def:b3:membrane-traffic:fusion}{SNARE} geknipt komt
er geen glycine of GABA vrij, vuren de motorische neuronen
ongeremd en trekt elke spier samen --- een krampachtige verlamming.
Enkele nanogrammen botulinetoxine in een overactieve spier gespoten
leggen die plaatselijk maandenlang stil (tot er nieuw \omterm{def:b3:membrane-traffic:fusion}{SNARE} is gemaakt
en het uiteinde herstelt): een behandeling voor dystonieën,
spasticiteit, scheelzien en rimpels.
\end{solution}

\begin{solution}{pb:b3:membrane-traffic:1}
\textbf{1.} $10^{7}\times \qty{25000}{Da}\times 1.66\times 10^{-24}$ g
$= \qty{0.42}{pg}$ per minuut, \qty{600}{pg} per dag --- tweemaal haar
eigen eiwitgehalte.
\textbf{2.} Volume van een blaasje $\tfrac{4}{3}\pi\,35^{3} = 1.8\times 10^{5}$
nm$^{3}$, waarvan \qty{30}{\%} $5.4\times 10^{4}$ is; één enzymmolecuul
$\tfrac{4}{3}\pi\,2^{3} = 34$ nm$^{3}$: ongeveer \num{1600} moleculen
per blaasje; $10^{7}/1600 \approx 6200$ blaasjes per minuut.
\textbf{3.} Oppervlak $4\pi\,35^{2} = \qty{0.0154}{\micro m^2}$ elk, \qty{95}{\micro
m^2} per minuut: de \qty{6000}{\micro m^2} van het reticulum in
ongeveer een uur als er niets terugkwam.
\textbf{4.} $\qty{95}{\micro m^2} = 9.5\times 10^{7}$ nm$^{2}$, twee lagen
bij \qty{0.7}{nm^2}: $2.7\times 10^{8}$ lipiden per minuut. De cel
brengt membraan terug met COPI-blaasjes en recycling, en maakt lipide
aan om te vervangen wat definitief in korrels vertrekt.
\textbf{5.} Volume van een korrel $5.2\times 10^{8}$ nm$^{3}$, waarvan \qty{30}{\%}
$1.6\times 10^{8}$: $4.7\times 10^{6}$ moleculen per korrel; $6\times
10^{8}$ moleculen per uur vullen ongeveer $130$ korrels.
\textbf{6.} Elke korrel voegt $\pi d^{2} = \qty{3.1}{\micro m^2}$ toe; $300$
voegen \qty{940}{\micro m^2} toe aan \qty{50}{\micro m^2}: twintigvoudig.
Het membraan moet binnen minuten door compenserende endocytose worden
teruggehaald.
\textbf{7.} $t^{*} = \ln 2/0.05 = \qty{13.9}{min}$, $B = 0.50$.
\textbf{8.} $A = e^{-3} = 0.05$; $B = 2(e^{-1.5} - e^{-3}) = 0.35$;
$C = 0.60$.
\textbf{9.} \qty{10}{min} en \qty{20}{min}; gemiddeld totaal
\qty{30}{min}.
\textbf{10.} $t^{*} = \ln(0.1)/(-0.09) = \qty{25.6}{min}$; $B(t^{*}) =
(0.1/(-0.09))(e^{-2.56} - e^{-0.256}) = 0.77$: het golgiapparaat zwelt
op, volgepropt met drie kwart van de vracht.
\textbf{11.} Bij het maximum geldt $k_{1}e^{-k_{1}t^{*}} = k_{2}e^{-k_{2}t^{*}}$;
invullen in $B$ geeft $B_{\max} = (k_{1}/k_{2})^{\,k_{2}/(k_{2} -
k_{1})}$: de macht is $k_{2}/(k_{2} - k_{1})$. Voor $k_{1} > k_{2}$ is
de exponent negatief en het grondtal boven $1$; voor $k_{1} < k_{2}$ is
het grondtal onder $1$ en de exponent positief --- in beide gevallen
ligt de waarde onder $1$ (controle: hier $2^{-1} = 0.5$).
\textbf{12.} Voor $k_{1} \gg k_{2}$ nadert $B(t) \to e^{-k_{2}t}$: het
merkstof zit meteen in het golgiapparaat en vertrekt met $k_{2}$ --- de
reticulumstap is onzichtbaar en de golgicurve is een eenvoudig verval.
\textbf{13.} $\theta = 2/(2+2) = 0.5$; $J = \num{50000}\times 0.2\times
0.1\times 0.5/(0.1 + 0.1) = 2500$ deeltjes per minuut, \num{150000} per
uur.
\textbf{14.} $1.5\times 10^{5}\times 1500 = 2.3\times 10^{8}$
cholesterolmoleculen per uur; $5\times 10^{9}/2.3\times 10^{8} \approx 22$ uur.
\textbf{15.} Deel aan het oppervlak $k_{r}/(k_{r} + k_{i}\theta) = 0.5$; één
rondreis kost $1/(k_{i}\theta) + 1/k_{r} = 10 + 10 = \qty{20}{min}$:
ongeveer $60$ reizen in twintig uur.
\textbf{16.} $J = 1250$ per minuut. Onder de verzadiging is $J \propto R\,[L]$,
dus moet het LDL in het plasma bij de helft van de receptoren
verdubbelen voor dezelfde klaring.
\textbf{17.} $R$ terug op \num{50000}: de klaring is hersteld, het LDL in
het plasma zakt weer naar normaal --- het effect van de statine.
\textbf{18.} Zonder receptoren wordt het LDL alleen langs trage
niet-specifieke wegen geklaard, zodat het dagen in plaats van uren
circuleert, wordt geoxideerd en wordt opgenomen door de
opruimreceptoren van macrofagen, die de schuimcellen van
atherosclerotische plaques worden; en statines, die werken door
receptoren aan te zetten, hebben niets om aan te zetten.
\textbf{19.} $V = \tfrac{4}{3}\pi(0.25)^{3} = \qty{0.065}{\micro m^3} =
6.5\times 10^{-17}$ L. Bij pH $7.2$: $6.3\times 10^{-8}\times 6.5\times
10^{-17} = 4\times 10^{-24}$ mol --- ongeveer $2$ vrije protonen. Bij pH $4.7$:
$2\times 10^{-5}\times 6.5\times 10^{-17} = 1.3\times 10^{-21}$ mol,
ongeveer $780$.
\textbf{20.} $50\times 780 \approx \num{39000}$ protonen; $50$ pompen die
$200$ per seconde verzetten halen \num{10000} per seconde: ongeveer
\qty{4}{s}.
\textbf{21.} Positieve lading naar binnen pompen maakt het lumen
positief, en enkele duizenden ladingen zouden de pomp stilleggen;
chloride komt door een kanaal naar binnen (of kationen gaan naar
buiten) om de lading te neutraliseren.
\textbf{22.} Enzymen die bij neutrale pH inactief zijn richten geen
schade aan als een \omterm{def:b3:membrane-traffic:endocytosis}{lysosoom} lekt of als zij verkeerd naar het cytosol
of het oppervlak worden gestuurd: de eis van zuur houdt de vertering
binnen het compartiment dat ervoor is gebouwd.
\textbf{23.} Beide moeten een pH-sensor dragen --- histidines, waarvan
de protonering rond pH~6 het bindingsoppervlak verandert: de
LDL-receptor vouwt bij protonering een domein over zijn eigen
bindingsplaats, en de mannose-6-fosfaatreceptor verliest evenzo haar
affiniteit.
\textbf{24.} $10^{7.2 - 4.7} = 10^{2.5} \approx 300$-voudige ophoping;
de gevangen base verbruikt protonen, verhoogt de pH van het \omterm{def:b3:membrane-traffic:endocytosis}{lysosoom} en
schakelt de hydrolasen uit --- zo doodt chloroquine malariaparasieten in
hun voedselvacuole en zo blokkeert het de \omterm{def:b3:membrane-traffic:endocytosis}{autofagie} in experimenten.
\textbf{25.} Het golgimaximum bij \qty{14}{min}; zo'n \num{6200} blaasjes
die per minuut het reticulum verlaten; \num{150000} LDL-deeltjes
opgenomen per uur.
\end{solution}
